% !TeX root = ./main.tex

% --------------------------------------------------
% 資訊設定（Information Configs）
% --------------------------------------------------

\ntusetup{
  university*   = {National Taiwan University},
  university    = {國立臺灣大學},
  college       = {電機資訊學院},
  college*      = {College of Electrical Engineering and Computer Science},
  institute     = {電子工程學研究所},
  institute*    = {Graduate Institute of Electronics Engineering}, % 請確認您的科系為Department of xxx或Institute of xxx
  title         = {具有容錯性的無歧義量子態鑑別},
  title*        = {Unambiguous Quantum State Discrimination with Error Tolerance},
  author        = {馬健凱},
  author*       = {Chien-Kai Ma},
  ID            = {R11943106},
  advisor       = {江介宏},
  advisor*      = {Jie-Hong Roland Jiang},
  date          = {2025-07-22},         % 若註解掉，則預設為當天
  oral-date     = {2025-07-22},         % 若註解掉，則預設為當天
  DOI           = {10.6342/NTU202502623},
  keywords      = {無歧義量子態鑑別, 量子電路, 正算子值測度, 去極化雜訊},
  keywords*     = {UQSD, quantum circuits, POVM, depolarizing noise},
}

% --------------------------------------------------
% 加載套件（Include Packages）
% --------------------------------------------------

\usepackage[numbers, sort&compress]{natbib}     % 參考文獻
\usepackage{amsmath, amsthm, amssymb}           % 數學環境
\usepackage{ulem, CJKulem}                      % 下劃線、雙下劃線與波浪紋效果
\usepackage{booktabs}                           % 改善表格設置
\usepackage{multirow}                           % 合併儲存格
\usepackage{diagbox}                            % 插入表格反斜線
\usepackage{array}                              % 調整表格高度
\usepackage{longtable}                          % 支援跨頁長表格
\usepackage{paralist}                           % 列表環境
\usepackage{quantikz}                           % Quantum
\usepackage{minted}                             % Code blocks


\usepackage{lipsum}                     % 英文亂字
\usepackage{zhlipsum}                   % 中文亂字

% --------------------------------------------------
% 套件設定（Packages Settings）
% --------------------------------------------------
\setCJKmonofont{Noto Sans Mono CJK TC}  % for Traditional Chinese
\setCJKsansfont{Noto Sans CJK TC}
